\documentclass[paper=screen,mode=present,style=zysimple]{powerdot}

\pdsetup{lf={Theory and Practice of Functional Programming}}

\title{{\small \textsc{Foundations of Programming Languages}}
\\[12pt]
{\Large \color{TitleColor} Theory and Practice of Functional Programming}
}
\author{\small \textbf{Dr. ZHANG Yu}
  \\[2pt] \small Institute of Software, Chinese Academy of Sciences 
  \\[5pt] \includegraphics[height=.04\slideheight]{iscas.eps} 
  \qquad \qquad 
  \includegraphics[height=.04\slideheight]{cas.eps}
}
\date{\scriptsize Fall term, 2011 \qquad GUCAS, Beijing}

\begin{document}

\maketitle 

%% \begin{slide}[toc=]{Outline}
%% \tableofcontents
%% \end{slide}

\begin{slide}{Theory and Practice of Functional Programming}
The course is officially entitled {\em Foundations of Programming Languages}, but since last year it focuses on 
\emblue{functional programming} --- not just as a theoretical model for PLs, but also as practical languages 
to build real applications.
\begin{itemize}
\item Functional languages are becoming popular in recent years and used to build real applications in industry:
\begin{itemize}
\item Erlang has been used for almost two decades in Ericsson; Twitter switched to scala a few years ago;  
  part of Amazon Web Services is written in Erlang ...
\item Finance industry starts to make intensive use of FPL: 
  Erlang in Goldman Sachs and Deutsche Bank, Ocaml in Jane Street, Scala in UBS Investment Bank, 
  F\# and Haskell in Credit Swiss, ...
\end{itemize}
\item Programming in FPL can help to understand abstract concepts in theory.
\end{itemize}
\end{slide}

\begin{slide}{Why Haskell?}
We shall use Haskell as the main language in the course (with a comparison to Ocaml).
\begin{itemize}
\item Haskell is known to be hard for non-FPL programmers.
\item Haskell is rich, covering most theoretical concepts of FP and type systems.
\item Haskell is a \emred{pure} functional language.
\end{itemize}
\end{slide}

\begin{slide}{Course outline}
10 lectures (3 hours each):
\begin{itemize}
\item \emblue{Foundation and theory} of FPL.
\begin{itemize}
\item Lambda-calculus, simple type theory, semantic model, modeling of imperative programs, ...
\item Basic category concepts, monads, concurrency, ...
\end{itemize}
\item \emblue{Practical programming} using FPL.
\begin{itemize}
\item Programming in Haskell: lists, type classes, module systems, functors, ...
\item Data structure and algorithms in FPL.
\end{itemize}
\end{itemize}
\end{slide}

\begin{slide}{Schedule}
\hspace*{4em}
\begin{tabular}[b]{l || l }
\hline \\[-0.8em]
Week 1 & Lambda calculus
\\[0.5em] \hline \\[-0.8em]
Week 2  & Simply typed lambda-calculi (PCF)
\\[0.5em] \hline \\[-0.8em]
Week 3 & Type theory 
\\[0.5em] \hline \\[-0.8em]
Week 4 & Introduction to Haskell 
\\[0.5em] \hline \\[-0.8em]
Week 5, 6 & Functional data structure and algorithms
\\[0.5em] \hline \\[-0.8em]
Weak 7, 8 & Advanced Haskell programming 
\\[0.5em] \hline \\[-0.8em]
Weak 9 & Advanced topics in type theory
\\[0.5em] \hline \\[-0.8em]
Week 10 & Course review
\\[0.5em] \hline 
\end{tabular}
\\
This is a rough schedule --- after the introduction of Haskell, theory and programming practice will be mixed.
\end{slide}

\begin{slide}{Reference}
\begin{itemize}
\item There is no textbook for the course. I list some reference books on the course page:
\begin{center}
{\em \empurple{http://lcs.ios.ac.cn/\~{}yzhang/fopl/}}
\end{center}
More references will be given during the course and available on the page, 
but you are also encouraged to do some research \emred{on your own}.
\\
Lecture notes and other material will  be download-able from the page.
\item Me: \\
I will be available for answering your questions {\em every Tuesday afternoon from 15:00 to 18:00}, 
at \emblue{Room 315 of ISCAS building NO. 5}.
\end{itemize}
\end{slide}

\begin{slide}[toc=]{Evaluation}
Evaluation will base on 
\begin{itemize}
\item {\bf Projects}: there will be several programming pearls during the course, 
  and one or two course projects on functional programming.
  
\item {\bf Final exam}: this consists of all the theoretical contents that are covered by the course, 
  as well as functional programming concepts. \\
  The exam will take place on \emred{Friday January 13, 2012}.
\end{itemize}
The two parts will be roughly half-half in your final score. 
\end{slide}

\end{document}
 

